% ---
% Reemplaza líneas 3395–3449 de Docs/Technical_Report.tex.
% Decisiones que lo respaldan: D-005, D-006, D-009, D-022, H-002, H-050.
% ---

\subsection{Hallazgos principales}

El análisis exploratorio permite identificar los
siguientes hallazgos de mayor relevancia para el
diseño del pipeline experimental:

\begin{enumerate}

% fuente: D-005, RF-19; 1_EDA.ipynb
\item \textbf{Desbalance de clases diferenciado.}
El ratio benigno/maligno es 1.16:1 para masas y
1.78:1 para calcificaciones. La ponderación de
clases durante el entrenamiento del MLP es
necesaria para calcificaciones pero opcional para
masas. Los modelos se evalúan con AUC-ROC como
métrica principal por su robustez ante este
desbalance.

% fuente: H-050
\item \textbf{Asimetría de separabilidad esperada.}
Las masas presentan patrones morfológicos
claramente diferenciados entre clases
(IRREGULAR/\allowbreak SPICULATED vs.\ OVAL/\allowbreak CIRCUMSCRIBED),
prediciendo mayor separabilidad en el espacio
de representación. Las calcificaciones comparten
el tipo morfológico más frecuente (PLEOMORPHIC)
entre ambas clases, prediciendo menor separabilidad
y constituyendo el caso de prueba más exigente
para los quantum embeddings.

% fuente: D-006
\item \textbf{Variables clínicas no predictivas.}
La densidad mamaria y el \textit{subtlety}
tienen correlación prácticamente nula con la
patología. La información discriminativa reside
en los patrones morfológicos de la imagen,
validando el uso de PyRadiomics como extractor
de características.

% fuente: D-009, H-051; 1_EDA.ipynb
\item \textbf{Imágenes de alta dimensionalidad.}
Las mamografías completas tienen dimensiones
medianas de 5\,037\,$\times$\,2\,960 píxeles,
incompatibles con la entrada directa al circuito
cuántico. El trabajo con recortes de la región de
interés a resolución original, seguido de
extracción de características radiómicas y selección
de $k$ características, es la estrategia viable para
la simulación cuántica disponible.

% fuente: D-022, H-002; 1_metadatos_dicom.csv
\item \textbf{Convención fotométrica uniforme.}
La totalidad del dataset utiliza \texttt{MONOCHROME2},
eliminando la necesidad de inversión de píxeles
en el preprocesamiento.

\end{enumerate}
